\documentclass[11pt,a4paper]{article}
\usepackage[utf8]{inputenc}
\usepackage[T1]{fontenc}
\usepackage[french]{babel}
\usepackage[margin=2.5cm]{geometry}
\usepackage{graphicx}
\usepackage{hyperref}
\usepackage{listings}
\usepackage{xcolor}
\usepackage{float}

\definecolor{codegreen}{rgb}{0,0.6,0}
\definecolor{codegray}{rgb}{0.5,0.5,0.5}
\definecolor{codepurple}{rgb}{0.58,0,0.82}
\definecolor{backcolour}{rgb}{0.95,0.95,0.92}

\lstdefinestyle{mystyle}{
    backgroundcolor=\color{backcolour},   
    commentstyle=\color{codegreen},
    keywordstyle=\color{magenta},
    numberstyle=\tiny\color{codegray},
    stringstyle=\color{codepurple},
    basicstyle=\ttfamily\footnotesize,
    breakatwhitespace=false,         
    breaklines=true,                 
    captionpos=b,                    
    keepspaces=true,                 
    numbers=left,                    
    numbersep=5pt,                  
    showspaces=false,                
    showstringspaces=false,
    showtabs=false,                  
    tabsize=2
}
\lstset{style=mystyle}

\title{\textbf{Projet C++ : Générateur de Donjon et Labyrinthe Interactif} \\ \large Rapport de Projet - Master 1 Sciences pour l'Ingénieur}
\author{BOUANZOUL Ahmed Walid \and DEBZ Wail}
\date{7 Mai 2026}

\begin{document}

\maketitle

\section{Introduction}
Ce projet a pour objectif la conception et l'implémentation d'un jeu d'exploration de donjon généré de manière procédurale, développé en C++17. Il met en pratique les concepts fondamentaux de la programmation orientée objet (POO) tels que l'héritage, le polymorphisme, et les patrons de conception, ainsi que l'utilisation d'algorithmes classiques pour la génération du labyrinthe et la recherche du chemin optimal.

\section{Architecture Globale et Classes Principales}
Le projet est organisé selon une architecture modulaire séparant les déclarations (\texttt{.h}) et les implémentations (\texttt{.cpp}).

\subsection{Hiérarchie des cases et Polymorphisme}
Le cœur de la modélisation repose sur la classe abstraite \texttt{Case}, qui définit l'interface commune de tous les éléments de la grille du donjon. Elle déclare des méthodes virtuelles pures telles que \texttt{afficher()} et \texttt{appliquerEffet(Aventurier\&)}, garantissant un comportement polymorphe.
Cinq classes concrètes héritent de \texttt{Case} :
\begin{itemize}
    \item \texttt{Mur} (\texttt{\#}) : obstacle infranchissable.
    \item \texttt{Passage} (\texttt{ }) : espace libre permettant le déplacement.
    \item \texttt{Tresor} (\texttt{+}) : permet au joueur d'augmenter son inventaire.
    \item \texttt{Monstre} (\texttt{M}) : déclenche un combat infligeant des dégâts ou offre la possibilité de fuir.
    \item \texttt{Piege} (\texttt{T}) : inflige automatiquement des dégâts à l'aventurier.
\end{itemize}

\subsection{Patron de conception \textit{Factory}}
Pour gérer l'instanciation de ces différentes cases sans coupler fortement le code de génération aux classes concrètes, nous avons implémenté le \textit{Design Pattern Factory}. La classe \texttt{CaseFactory} expose une méthode statique \texttt{creerCase(TypeCase type)} qui retourne un pointeur \texttt{Case*}. Ce choix architectural respecte le principe d'ouverture/fermeture (Open/Closed) : l'ajout d'un nouveau type de case ne nécessite de modifier que l'énumération \texttt{TypeCase} et la Factory, sans impacter la logique de génération du \texttt{Donjon}.

\subsection{Les classes Donjon et Aventurier}
\begin{itemize}
    \item \texttt{Donjon} : Gère la grille bidimensionnelle (\texttt{std::vector<std::vector<Case*>>}). Elle orchestre la génération de la carte, l'affichage et l'algorithme de calcul de chemin. La gestion de la mémoire y est rigoureuse : les destructeurs libèrent proprement les pointeurs alloués par la Factory pour éviter toute fuite de mémoire.
    \item \texttt{Aventurier} : Encapsule l'état du joueur (position, points de vie, inventaire) et contient la boucle de jeu principale. Elle interagit avec le donjon via la méthode \texttt{resoudreCase(Case*)}.
\end{itemize}

\section{Algorithmes Utilisés : Choix et Implémentation}
\subsection{Génération du Donjon : Recursive Backtracking}
\textbf{Choix et Justification :} Pour garantir la création d'un ``labyrinthe parfait'' (un graphe connexe acyclique où il n'existe qu'un unique chemin entre deux points), nous avons opté pour l'algorithme de \textit{Recursive Backtracking} (Parcours en profondeur avec retour arrière). Ce choix assure que chaque zone de la grille est accessible.
\\
\textbf{Implémentation :} L'algorithme part d'une case initiale (1, 1), sélectionne aléatoirement une direction parmi les 4 cardinales, vérifie si la case à une distance de 2 est non visitée (constituée de murs). Si c'est le cas, il ``casse'' le mur intermédiaire en transformant les deux cases en \texttt{Passage}, et s'appelle récursivement. Pour un choix équitable et aléatoire des directions, nous avons utilisé un \texttt{std::mt19937} (Mersenne Twister) couplé à \texttt{std::shuffle}. Les entités (Trésors, Monstres, Pièges) sont ensuite placées aléatoirement sur les cases de type \texttt{Passage} selon des probabilités définies.

\subsection{Recherche de chemin : Algorithme BFS}
\textbf{Choix et Justification :} L'algorithme BFS (\textit{Breadth-First Search}) a été sélectionné pour calculer le chemin optimal vers la sortie. La grille de notre donjon n'étant pas pondérée (chaque déplacement coûte 1), l'algorithme BFS garantit de trouver le chemin le plus court. De plus, sa complexité de $\mathcal{O}(V + E)$ est parfaitement adaptée aux dimensions de notre carte.
\\
\textbf{Implémentation et Structures de données :} Nous avons utilisé la file \texttt{std::queue<std::pair<int, int>>} de la bibliothèque standard (STL) pour explorer les voisins niveau par niveau. Deux tableaux 2D (\texttt{std::vector}) gardent une trace des cases visitées (\texttt{visite}) et des cases parentes (\texttt{parent}) pour reconstruire le chemin en remontant de la sortie vers le joueur.

\section{Difficultés Rencontrées et Solutions}
\begin{itemize}
    \item \textbf{Gestion dynamique de la mémoire et fuites (Memory Leaks) :} La manipulation de pointeurs de la classe de base \texttt{Case*} pointant vers des objets dérivés a initialement causé des fuites de mémoire. \\
    \textit{Solution :} Nous avons implémenté un destructeur virtuel (\texttt{virtual \char`\~Case() = default;}) dans la classe mère et nous nous sommes assurés que le destructeur du \texttt{Donjon} itère sur la grille complète pour appeler \texttt{delete} sur chaque instance. Nous avons également corrigé les fuites lors du remplacement des cases pendant le placement des entités en supprimant d'abord l'ancienne case avant d'assigner la nouvelle.
    \item \textbf{Problème d'inclusion circulaire :} Les classes \texttt{Aventurier} et \texttt{Case} interagissaient mutuellement, causant des erreurs de compilation dues aux inclusions circulaires (notamment pour la méthode \texttt{appliquerEffet}). \\
    \textit{Solution :} Nous avons eu recours aux \textit{forward declarations} (déclarations anticipées) comme \texttt{class Aventurier;} dans les fichiers d'en-tête, limitant l'inclusion des fichiers \texttt{.h} aux fichiers sources \texttt{.cpp}.
    \item \textbf{Scintillement de la console :} Un affichage caractère par caractère dans la boucle de jeu provoquait un scintillement (flickering) désagréable. \\
    \textit{Solution :} Nous avons optimisé l'affichage en construisant la grille entière dans un tampon (\texttt{std::ostringstream}) avant de l'afficher en une seule fois via \texttt{std::cout}.
\end{itemize}

\section{Exemples de Donjons Générés}
Afin de prouver la robustesse de notre génération procédurale, notre programme sauvegarde et génère des grilles. Voici la représentation ASCII d'un labyrinthe typique obtenu (identique au format des fichiers \texttt{.txt} de sauvegarde) :
\begin{verbatim}
+---------------------+
|#@#      #           #|
|# #+### ### ####### #|
|#  #   #     #   #M #|
|# ### ### # ### # # #|
|#     #   # #   # + #|
|##### # ### # ### # #|
|#   # #   # #     # #|
|# # # ### # ####### #|
|# #   #   # #       #|
|# ##### ### # ##### #|
|# #   #   # #     # #|
|# # # ### ####### # #|
|#   #     #   +   # #|
|###T###### # # ### #|
|#         # # # M   #|
|# ####### # ##### # #|
|# #     # #       # #|
|# # ### ########### #|
|#   # M        T  @E#|
+---------------------+
Position:(1,1) HP:100 Inv:0
\end{verbatim}
Le symbole \texttt{@} indique la position de départ (ou courante) du joueur, \texttt{E} la sortie, \texttt{\#} les murs, \texttt{+} les trésors, \texttt{M} les monstres, et \texttt{T} les pièges. 

\section{Journal de Prompts}
Afin d'accélérer notre développement et d'améliorer la qualité du code, nous avons utilisé des outils d'Intelligence Artificielle générative comme assistants de programmation pour le débogage et l'explication de concepts. Voici un extrait de notre journal de prompts :
\begin{enumerate}
    \item \textit{``Comment éviter le scintillement (flickering) lors du rafraîchissement d'une carte ASCII en C++ dans le terminal ?''} \\
    $\rightarrow$ L'outil nous a orientés vers l'utilisation du buffering avec \texttt{std::ostringstream}.
    \item \textit{``Quelle est la différence entre enum et enum class en C++ moderne ?''} \\
    $\rightarrow$ L'assistant a détaillé la sécurité de typage accrue de \texttt{enum class}, évitant les conversions implicites en entier. Nous l'avons appliqué pour \texttt{TypeCase}.
    \item \textit{``Mon destructeur de Donjon ne supprime pas correctement les objets enfants comme Tresor et Monstre à travers le pointeur Case*, pourquoi ?''} \\
    $\rightarrow$ Réponse cruciale nous rappelant la nécessité du destructeur virtuel (\texttt{virtual \char`\~Case() = default;}) pour éviter les fuites de mémoire (comportement indéfini).
    \item \textit{``Comment bien structurer l'algorithme BFS pour reconstruire le chemin vers la sortie dans une grille représentée par un vector2D ?''} \\
    $\rightarrow$ L'explication nous a aidés à organiser les tableaux \texttt{visite} et \texttt{parent}, permettant d'éviter les boucles infinies et de reconstruire le chemin efficacement.
\end{enumerate}

\section{Conclusion}
Ce projet a été l'occasion de consolider nos acquis en programmation C++. Le respect de la programmation orientée objet, combiné à l'utilisation du patron de conception Factory et de structures de données adaptées pour les algorithmes (BFS, Backtracking), nous a permis de produire un code maintenable, robuste et d'atteindre l'ensemble des objectifs fixés.

\end{document}
